    \section{Simulation Experiments}
    \label{section: simulation}
    In addition to the experiment conducted on the real embedded board above, we also designed simulation experiments to evaluate the system performance more comprehensively. Different from the real-world experiments, the simulation experiment results are obtained based on the average and the distribution of a large number of random task sets to cover a wider range of scenarios.
    
    \subsection{Simulated Task Sets Generation}

    \agent{notice that our simulation expriments use the precise RTA estimation, rather than the pessimistic RTA estimation from analytical derivation.}

    In simulation, we assume the robot is operated in a two-dimensional environment where we can use the \textit{x} and \textit{y} coordinates to describe a robot's location.
    For simplicity, the 2D environment is limited to a square of width $W^E$ and height $H^E$; in our experiments $W^E = H^E = 200$ m and the robot moves at $1.0$ m/s.
    During operation, the robot needs to run several computational tasks to ensure both safety and functionality. These tasks in our simulation experiments are generated randomly to cover different scenarios. The task set generation follows the typical real-time systems' experiments design~\cite{Wang2023RTAS, Zhao2020AnOF} based on Uunifast~\cite{Davis2008EfficientES}.
    Each task set comprises $N$ tasks scheduled on multiple ($2$ in experiments) CPU cores under partitioned scheduling, where $N \in \{4, 6, 8, 10, 12, 14, 16\}$ to cover a range of system scales; for each $N$ we generate $10$ independent task sets, and all reported results are aggregated over these sets.
    The tasks' periods are drawn from a limited set based on real-world systems~\cite{Percept21_RTSS, Kramer15benchmark}: $\{20, 33, 50, 100, 200, 500, 1000\}$ ms.
    The per-core CPU utilization is sampled uniformly from $[0.5, 1.0]$ and the Uunifast algorithm distributes this total across the $N$ tasks to set each task's \textit{average} execution time $\mu^E_{c,i}$.
    All tasks use implicit deadlines, $D_i = T_i$.

    % TODO: no need to re-introduce these 2 tasks' concepts, simplify this paragraph
    To capture the heterogeneity of real robotic workloads, each task is one of the two types defined in Section~\ref{section: system_model_tasks}:
    \textbf{Environment-dependent tasks with spatial variation} (the non-degenerate case of $\tau^E_i$) are $N_{\text{env}} = \lceil \rho_{\text{env}} \cdot N \rceil$ tasks, where $\rho_{\text{env}} \sim U[0.1, 0.3]$ is sampled per task set and the count is clamped to at least one; these tasks are drawn with probability proportional to their utilization and each is subject to a per-task utilization cap of $0.27$. Their execution times use the full GMM model with nonzero spatial correlations, described below.
    Each \textbf{QoS-configurable task} ($\tau^Q_i$) is formed by taking each remaining task independently with probability $0.5$; it has near-deterministic execution times ($\sigma_{c,i} = 0.001\,\mu^E_{c,i}$) and a discrete QoS-budget grid of $10$ options spanning $[0.05, 0.9] \times T_i$.
    The \textbf{remaining environment-dependent tasks} are the degenerate case of $\tau^E_i$ (Definition~\ref{def_env_task}), following the GMM with zero spatial correlation ($\rho_{x,c,i} = \rho_{y,c,i} = 0$).

    We used Gaussian Mixture Model (GMM) to simulate the execution time distribution's dependency on the environment. Specifically, we assume that the robot's physical location $(x, y)$ and each task's execution time distribution $\extDist{i}$ follows a joint 3D Gaussian distribution; the per-task distribution is a mixture of $K{=}4$ such components with mixture weights (each floored at $0.1$). Each component's probability density function $f^{E}()$ is
    \begin{equation}
        f^{E}(x, y, \extDist{i}) =\mathcal{N}(\boldsymbol{\mu}^E, \boldsymbol{\Sigma}^E)
        \label{eq: pdf_gmm}
    \end{equation}
    \begin{equation}
        \boldsymbol{\mu}^E = \begin{bmatrix}
            \mu^E_x & \mu^E_y & \mu^E_{c,i}
        \end{bmatrix}
    \end{equation}

\begin{equation}
    \boldsymbol{\Sigma}^E =
    \begin{bmatrix}
        \sigma_x^2 & 0 & \rho_{x,c,i} \sigma_x \sigma_{c,i} \\
        0 & \sigma_y^2 & \rho_{y,c,i} \sigma_y \sigma_{c,i} \\
        \rho_{x,c,i} \sigma_x \sigma_{c,i} & \rho_{y,c,i} \sigma_y \sigma_{c,i} & \sigma_{c,i}^2
    \end{bmatrix}
\end{equation}

    The values of the mean vector $\boldsymbol{\mu}^E$ and the covariance matrix $ \boldsymbol{\Sigma}^E$ in the probability density function~\eqref{eq: pdf_gmm} are decided as follows.
    Firstly, notice that the robot is assumed to be operated in a 2D square of width $W^E$ and height $H^E$. Therefore, by placing the coordinate origin at the center of the square, and assuming the probability density function covers 95\% confidence interval, and that gives us:
    \begin{equation}
        \mu^E_x = 0, \sigma_x = W^E/4
    \end{equation}
    \begin{equation}
        \mu^E_y = 0, \sigma_y = H^E/4
    \end{equation}
    As for the execution time $\extDist{i}$ of the task $\tau_i$, its average execution time $\mu^E_{c,i}$ is generated from the Uunifast algorithm, and its standard deviation $\sigma_{c,i}$ is a random portion of its average execution time: $\sigma_{c,i}=k\,\mu^E_{c,i}$ with $k \sim U[0.3, 0.4]$. For QoS-configurable tasks the ratio is instead fixed at $k = 0.001$, giving near-deterministic execution times.
    In the covariance matrix $\boldsymbol{\Sigma}^E$, the covariance ratios $\rho_{x,c,i}, \rho_{y,c,i}$ capture the execution time distribution $\extDist{i}$'s dependency on the physical environment: they are nonzero only for environment-dependent tasks with spatial variation, where $\rho_{x,c,i} \sim U[-0.9, -0.7]$ and $\rho_{y,c,i} \sim U[-0.1, 0.1]$ are sampled per GMM component;
    for constant-distribution environment-dependent tasks and QoS-configurable tasks both correlations are zero.
    The marginal execution-time bounds of any environment-dependent task are set to $\mu^E_{c,i} \pm 2\sigma_{c,i}$, while a QoS-configurable task's bounds span its QoS-budget grid.

    Each task $\tau_i$ also carries two SP-specific parameters sampled per task: its importance weight $w_i \sim U[0.1, 1.0]$ (subsequently normalized so that $\sum_i w_i = 1$), and its deadline-miss probability threshold $\Theta_i \sim U[0.001, 0.9]$, which is the probability of missing its deadline above which the task no longer contributes its full safety value.
    The important-task subset $\boldsymbol{\tau}^{safe}$ is then taken as the top $50\%$ of tasks by weight, $\lceil 0.5 N \rceil$ tasks, modeling safety-critical tasks that must be shielded from overload, while the remaining soft tasks may be shed when the system is under high computation load.
    Tasks are assigned to cores by a worst-fit (longest-utilization-first) heuristic to balance load across the two cores.

    \subsection{Compared Methods and Runtime Configuration}
    \label{section: sim_compared_methods}
    Each simulation runs for $600$~s with the scheduler re-optimized every $10$~s (the re-optimization interval), and each scheduler activation is bounded by a $1$~s optimization budget; the base random seed is $1000$ and is advanced per task set.
    We compare the following scheduling methods. The proposed method is \textbf{INCR}, the incremental optimizer with periodic re-optimization (the \textbf{INCR\_Reopt\_10} arm re-optimizes from scratch every $10$ intervals and otherwise optimizes incrementally).
    Against it we evaluate four baselines: \textbf{BF}, the brute-force priority-and-QoS optimizer used as the offline optimality reference; \textbf{DM\_FAST} and \textbf{DM\_SLOW}, deadline-monotonic priority assignments paired with, respectively, the smallest and largest QoS budget on each task's grid; and \textbf{CFS}, the Linux completely-fair scheduler baseline that uses no priority or QoS optimization.
    We further run ablation arms that isolate individual mechanisms of the proposed optimizer: \textbf{INCR\_NO\_TL} disables QoS-budget optimization; \textbf{INCR\_WCET} replaces the execution-time distribution with its pessimistic worst-case value; \textbf{INCR\_NO\_REOPT} disables periodic re-optimization (pure incremental); and \textbf{INCR\_Reopt\_1/30/60} vary the re-optimization period to study its effect.

    \subsection{Results}
    \agent{The figures and aggregation counts below are stale (they reflect a prior 100-taskset / 5-scheduler run: CFS, RM\_Fast, RM\_Slow, Brute-Force, Incremental). They must be regenerated from a prod run of the current configuration (10 tasksets per $N$, $N\in\{4,6,8,10,12,14,16\}$, schedulers INCR\_Reopt\_10 / BF / DM\_FAST / DM\_SLOW / CFS + ablations) before submission; the caption text and per-scheduler labels will change accordingly.}
	% \begin{figure*}[ht!]
	% 	\begin{subfigure}[c]{0.32\textwidth}
	% 		\centering
	% 		\includegraphics[width=\textwidth]{pictures/simulation/cpu_util_0_wolegend.pdf}
	% 		\caption{CPU 0 Utilization}
	% 		\label{fig_sim_cpu0}
	% 	\end{subfigure}
	% 	\begin{subfigure}[c]{0.32\textwidth}
	% 		\centering
	% 		\includegraphics[width=\textwidth]{pictures/simulation/cpu_util_1_wolegend.pdf}
	% 		\caption{CPU 1 Utilization}
	% 		\label{fig_sim_cpu}
	% 	\end{subfigure}
    
		
	% 	\caption{``Expected'' CPU Utilization variation of 100 random tasks during simulation experiments (See more in Section~\ref{section: simulation_exp}.} 
	% 	\label{figs_cpu_sim}
	% \end{figure*}

\begin{figure*}[ht!]
    \centering
    \begin{subfigure}[c]{0.32\textwidth}
        \centering
        \includegraphics[width=\textwidth]{pictures/simulation/cpu_util_0_wolegend.pdf}
        \caption{CPU 0 Utilization}
        \label{fig_sim_cpu0}
    \end{subfigure}
    \begin{subfigure}[c]{0.32\textwidth}
        \centering
        \includegraphics[width=\textwidth]{pictures/simulation/cpu_util_1_wolegend.pdf}
        \caption{CPU 1 Utilization}
        \label{fig_sim_cpu1}
    \end{subfigure}
    \caption{``Expected'' CPU Utilization variation of 100 random tasks during simulation experiments (See more in Section~\ref{section: simulation_exp}).} 
    \label{figs_cpu_sim}
\end{figure*}

    
	\begin{figure*}[ht!]
		\centering 
		\begin{subfigure}[t]{0.32\textwidth} 
			\centering
			\includegraphics[width=\textwidth]{pictures/simulation/box_plot_of_all_data_CFS.pdf}
			\caption{CFS}
			\label{fig_cfs_perf_sim}
		\end{subfigure}
		\begin{subfigure}[t]{0.32\textwidth}
			\centering
			\includegraphics[width=\textwidth]{pictures/simulation/box_plot_of_all_data_RM_FAST.pdf}
			\caption{RM\_Fast}
			\label{fig_rm_fast_sim}
		\end{subfigure}
		\begin{subfigure}[t]{0.32\textwidth} 
			\centering
			\includegraphics[width=\textwidth]{pictures/simulation/box_plot_of_all_data_RM_SLOW.pdf}
			\caption{RM\_Slow}
			\label{fig_rm_slow_sim}
		\end{subfigure}
		
		\begin{subfigure}[c]{0.32\textwidth}
			\centering
			\includegraphics[width=\textwidth]{pictures/simulation/box_plot_of_all_data_BR.pdf}
			\caption{Brute-Force Scheduler}
			\label{fig_opt_bf_sim}
		\end{subfigure}
		\begin{subfigure}[c]{0.32\textwidth}
			\centering
			\includegraphics[width=\textwidth]{pictures/simulation/box_plot_of_all_data_INCR.pdf}
			\caption{Incremental Scheduler}
			\label{fig_opt_incre_sim}
		\end{subfigure}
		
		\caption{The SP metric of different scheduling methods during the simulation experiments} 
		\label{figs_exp_results_all_sim}
	\end{figure*}
	
        